% Source: https://artofproblemsolving.com/community/c3h1443813_somewhat_similar_to_2014_mathcounts_national_sprint_24
% Stable identity: aops-topic-1443813-q01
\textbf{Anna.} She knows $B+C=14-A$. She can be sure that $B\neq C$ only if $14-A$ is odd (if it is even, $B=C=\dfrac{14-A}{2}$ is possible). So $A$ is odd.

\textbf{Brett.} From $B$ alone he must exclude $A=B$, $B=C$ and $A=C$. If $B\le6$ then $(A,C)=(B,14-2B)$ is possible, so $B\ge7$. If $14-B$ is even then $A=C=\dfrac{14-B}{2}$ is possible, so $B$ is odd. Conversely if $B\ge7$ is odd, then $A=B$ or $C=B$ would force the third number to be $14-2B\le0$, and $A+C=14-B$ is odd. Also $A+C\ge2$ gives $B\le12$. So $B\in\{7,9,11\}$.

\textbf{Chris.} Everyone now knows that $A$ is odd and $B\in\{7,9,11\}$. The possible triples $(A,B,C)$ are
\[
(1,7,6),\ (1,9,4),\ (1,11,2),\ (3,7,4),\ (3,9,2),\ (5,7,2).
\]
Chris knows $C$ and can determine everything only if his value of $C$ occurs once in this list. The values $C=4$ and $C=2$ occur several times, while $C=6$ occurs once. So $(A,B,C)=(1,7,6)$ and
\[
ABC=\boxed{42}.
\]
